%%%%%%%%%%%%%%%%%%%%%%%%%%%%%%
%Theorem
\newcounter{theo}[section] \setcounter{theo}{0}
\renewcommand{\thetheo}{\arabic{section}.\arabic{theo}}
\newenvironment{theo}[2][]{%
  \refstepcounter{theo}%
  \ifstrempty{#1}%
  {\mdfsetup{%
      frametitle={%
          \tikz[baseline=(current bounding box.east),outer sep=0pt]
          \node[anchor=east,rectangle,fill=blue!50]
          {\strut Satz~\thetheo};}}
  }%
  {\mdfsetup{%
      frametitle={%
          \tikz[baseline=(current bounding box.east),outer sep=0pt]
          \node[anchor=east,rectangle,fill=blue!50]
          {\strut Satz~\thetheo:~#1};}}%
  }%
  \mdfsetup{innertopmargin=10pt,linecolor=blue!50,%
    linewidth=2pt,topline=true,%
    frametitleaboveskip=\dimexpr-\ht\strutbox\relax
  }
  \begin{mdframed}[]\relax
    \label{#2}}
    {\nopagebreak\end{mdframed}}
%%%%%%%%%%%%%%%%%%%%%%%%%%%%%%
%Lemma
\newcounter{lem}[section]\setcounter{lem}{0}
\renewcommand{\thelem}{\arabic{section}.\arabic{lem}}
\newenvironment{lem}[2][]{%
  \refstepcounter{lem}%
  \ifstrempty{#1}%
  {\mdfsetup{%
      frametitle={%
          \tikz[baseline=(current bounding box.east),outer sep=0pt]
          \node[anchor=east,rectangle,fill=green!20]
          {\strut Lemma~\thelem};}}
  }%
  {\mdfsetup{%
      frametitle={%
          \tikz[baseline=(current bounding box.east),outer sep=0pt]
          \node[anchor=east,rectangle,fill=green!20]
          {\strut Lemma~\thetheo:~#1};}}%
  }%
  \mdfsetup{innertopmargin=10pt,linecolor=green!20,%
    linewidth=2pt,topline=true,%
    frametitleaboveskip=\dimexpr-\ht\strutbox\relax
  }
  \begin{mdframed}[]\relax%
    \label{#2}}{\end{mdframed}}

%%%%%%%%%%%%%%%%%%%%%%%%%%%%%%
%Definition
\newcounter{deff}[section]\setcounter{deff}{0}
\renewcommand{\thedeff}{\arabic{section}.\arabic{deff}}
\newenvironment{deff}[2][]{%
  \refstepcounter{deff}%
  \ifstrempty{#1}%
  {\mdfsetup{%
      frametitle={%
          \tikz[baseline=(current bounding box.east),outer sep=0pt]
          \node[anchor=east,rectangle,fill=red!60]
          {\strut Definition~\thedeff};}}
  }%
  {\mdfsetup{%
      frametitle={%
          \tikz[baseline=(current bounding box.east),outer sep=0pt]
          \node[anchor=east,rectangle,fill=red!60]
          {\strut Definition~\thedeff:~#1};}}%
  }%
  \mdfsetup{innertopmargin=10pt,linecolor=red!60,%
    linewidth=2pt,topline=true,%
    frametitleaboveskip=\dimexpr-\ht\strutbox\relax
  }
  \begin{mdframed}[]\relax%
    \label{#2}}{\end{mdframed}}

%%%%%%%%%%%%%%%%%%%%%%%%%%%%%%
% Beweisumgebung anpassen:
% Proof-Umgebung neu definieren mit fettem Titel
\makeatletter
\renewenvironment{proof}[1][Beweis]{%
  \par\pushQED{\qed}%
  \normalfont\topsep6pt \trivlist
  \item[\hskip\labelsep \bfseries #1.]\ignorespaces
}{%
  \popQED\endtrivlist\@endpefalse
}
\makeatother
\renewcommand{\qedsymbol}{$\blacksquare$}       % Ausgefülltes Quadrat

%%%%%%%%%%%%%%%%%%%%%%%%%%%%%%
%Example
\newcounter{examplebox}[section]\setcounter{examplebox}{0}
\renewcommand{\theexamplebox}{\arabic{section}.\arabic{examplebox}}
\newenvironment{examplebox}[2][]{%
  \refstepcounter{examplebox}%
  \ifstrempty{#1}%
  {\mdfsetup{%
      frametitle={%
          \tikz[baseline=(current bounding box.east),outer sep=0pt]
          \node[anchor=east,rectangle,fill=gray!20]
          {\strut Beispiel~\theexamplebox};}}
  }%
  {\mdfsetup{%
      frametitle={%
          \tikz[baseline=(current bounding box.east),outer sep=0pt]
          \node[anchor=east,rectangle,fill=gray!20]
          {\strut Beispiel~\theexamplebox:~#1};}}%
  }%
  \mdfsetup{innertopmargin=10pt,linecolor=gray!20,%
    linewidth=2pt,topline=true,%
    frametitleaboveskip=\dimexpr-\ht\strutbox\relax
  }
  \begin{mdframed}[]\relax%
    \label{#2}}
    {\end{mdframed}}
%%%%%%%%%%%%%%%%%%%%%%%%%%%%%%%%
%%%%%%%%%%%%%%%%%%%%%%%%%%%%%%
%Aufgabe
\newcounter{aufgabe}[section]\setcounter{aufgabe}{0}
\renewcommand{\theaufgabe}{\arabic{section}.\arabic{aufgabe}}
\newenvironment{aufgabe}{%
  \refstepcounter{aufgabe}%
  \vspace{0.5cm}
  \strut \textbf{Aufgabe~\theaufgabe}\\}
{\vspace{2cm}\\}
%%%%%%%%%%%%%%%%%%%%%%%%%%%%%%%
%%%%%%%%%%%%%%%%%%%%%%%%%%%%%%
%Fazit
\newenvironment{fazit}[2][]{%
  \ifstrempty{#1}%
  {\mdfsetup{%
      frametitle={%
          \tikz[baseline=(current bounding box.east),outer sep=0pt]
          \node[anchor=east,rectangle,fill=gray!20]
          {\strut Merke};}}
  }%
  {\mdfsetup{%
      frametitle={%
          \tikz[baseline=(current bounding box.east),outer sep=0pt]
          \node[anchor=east,rectangle,fill=gray!20]
          {\strut #1};}}%
  }%
  \mdfsetup{innertopmargin=10pt,linecolor=gray!20,%
    linewidth=2pt,topline=true,%
    frametitleaboveskip=\dimexpr-\ht\strutbox\relax
  }
  \begin{minipage}{\textwidth}
    \begin{mdframed}[]\relax%
      \label{#2}}{\end{mdframed}
  \end{minipage}}
%%%%%%%%%%%%%%%%%%%%%%%%%%%%%%%%
%%%%%%%%%%%%%%%%%%%%%%%%%%%%%%
%Beispiel in Absatz
\newcounter{example}[section]\setcounter{example}{0}
\renewcommand{\theexample}{\arabic{section}.\arabic{example}}

\newenvironment{example}{%
  \par\addvspace{0.75\baselineskip}% Mindestabstand oberhalb
  \noindent
  \begin{minipage}{\linewidth}% verhindert Seitenumbruch
    \refstepcounter{example}%
    \strut \textbf{Beispiel~\theexample}\\
    }{%
  \end{minipage}%
  \par\addvspace{0.75\baselineskip}% Mindestabstand unterhalb
}
%%%%%%%%%%%%%%%%%%%%%%%%%%%%%%%

\makeatletter
\newenvironment{chapquote}[2][2em]
{\setlength{\@tempdima}{#1}%
  \def\chapquote@author{#2}%
  \parshape 1 \@tempdima \dimexpr\textwidth-2\@tempdima\relax%
  \itshape}
{\par\normalfont\hfill--\ \chapquote@author\hspace*{\@tempdima}\par\bigskip}
\makeatother

\newcommand\kariert[2][0.5cm]{% 
  \begin{center}
    \begin{tikzpicture}[gray,step=#1]
      \pgfmathtruncatemacro\anzahl{(\linewidth-\pgflinewidth)/#1} % maximale Anzahl Kästchen pro Zeile
      \draw (0,0) rectangle (\anzahl*#1,#2*#1) (0,0) grid (\anzahl*#1,#2*#1);
    \end{tikzpicture}
  \end{center}
}
\newcommand{\R}{\mathbb{R}}
\newcommand{\Q}{\mathbb{Q}}
\newcommand{\N}{\mathbb{N}}